\documentclass[11pt,a4paper]{article}
\usepackage[T1]{fontenc}
\usepackage[utf8]{inputenc}
\usepackage[spanish,es-nodecimaldot]{babel}
\usepackage{lmodern}
\usepackage{geometry}
\usepackage{amsmath,amssymb}
\usepackage{booktabs,tabularx,array}
\usepackage{tikz}
\usetikzlibrary{arrows.meta,calc,patterns}
\usepackage{xcolor}
\usepackage{hyperref}
\usepackage{fancyhdr}
\usepackage{float}
\geometry{left=24mm,right=24mm,top=23mm,bottom=25mm}
\definecolor{armblue}{HTML}{176A93}
\definecolor{armred}{HTML}{BD3F3F}
\definecolor{armgray}{HTML}{4B5563}
\hypersetup{colorlinks=true,linkcolor=armblue,urlcolor=armblue}
\pagestyle{fancy}
\fancyhf{}
\fancyhead[L]{\small ARM-Structural | Banco de verificación 2D}
\fancyhead[R]{\small 25-09-2026}
\fancyfoot[L]{\small Revisión técnica independiente}
\fancyfoot[R]{\small \thepage}
\setlength{\headheight}{14pt}
\setlength{\parskip}{5pt}
\setlength{\parindent}{0pt}
\newcommand{\kN}{\,\mathrm{kN}}
\newcommand{\kNm}{\,\mathrm{kN\,m}}
\newcommand{\mm}{\,\mathrm{mm}}
\newcommand{\m}{\,\mathrm{m}}
\newcommand{\MPa}{\,\mathrm{MPa}}
\newcommand{\caseheading}[2]{\section{#1: #2}}
\newcommand{\note}[1]{\begin{quote}\small\color{armgray}#1\end{quote}}
\tikzset{member/.style={line width=1.25pt,draw=black},
         force/.style={-{Latex[length=2.5mm]},line width=1pt,draw=armred},
         axis/.style={-{Latex[length=2mm]},draw=armgray},
         defshape/.style={draw=armblue,dashed,line width=1.1pt}}

\begin{document}
\begin{titlepage}
\centering
\vspace*{22mm}
{\large ARM-Structural / Structural\_Kernel\par}
\vspace{13mm}
{\Huge\bfseries Banco de verificación\\[3mm]estructural 2D\par}
\vspace{7mm}
{\Large Cinco casos con referencias independientes, resultados y diagramas\par}
\vspace{19mm}
\begin{tabular}{@{}ll@{}}
Tipo & Informe técnico para contraste por ingenieros revisores\\
Versión & 1.0, 25 de septiembre de 2026\\
Modelo & Pórtico plano, elasticidad lineal, elemento Timoshenko\\
Unidades & kN, m, kN\,m, mm, mrad\\
Motores & TypeScript y Rust/WASM\\
Estado & Evidencia de verificación parcial; no certificación\\
\end{tabular}
\vfill
\begin{minipage}{0.86\textwidth}
\color{armgray}
Este informe permite repetir cálculos y detectar discrepancias. Los casos no
validan por sí solos análisis sísmico, diseño por resistencia, detallado de
armaduras ni cumplimiento del RNE. La aprobación del software requiere
ensayos adicionales, revisión de hipótesis, control de versiones y dictamen
firmado por profesionales competentes.
\end{minipage}
\end{titlepage}

\tableofcontents
\newpage

\section{Alcance, método y convenciones}
Se comparan cinco modelos planos de barras prismáticas con pequeñas
deformaciones, material elástico lineal y deformación por flexión y cortante.
Se anula el peso propio para aislar las cargas indicadas. No se aplican
factores normativos: los multiplicadores de carga muerta, viva y nodal son
unitarios. El eje global $X$ apunta a la derecha, $Y$ hacia arriba; $u_x$ y
$u_y$ son positivos en dichos sentidos. En las tablas de desplazamientos,
$\theta$ se expresa con la convención de la interfaz ARM, positiva en sentido
horario. Las referencias de diagramas emplean momento sagante positivo y
cortante positivo en la cara izquierda. Los momentos de reacción que muestra
ARM usan signo opuesto al momento físico antihorario; cuando importa se
informa la magnitud y se explica el equilibrio.

Los casos EX-01 a EX-04 tienen referencias cerradas de estática y teoría de
Timoshenko. EX-05 tiene una referencia separada por compatibilidad y método
de fuerzas (archivo de prueba
\texttt{tests/fixtures/portal-force-reference.mjs}); no es una medición
experimental ni una solución cerrada publicada. Ambos motores, TypeScript
($\mathrm{TS}$) y Rust/WASM ($\mathrm{R}$), se ejecutaron sobre el mismo modelo.
La concordancia TS--Rust es una comprobación de paridad, no sustituye la
referencia independiente.

Para sección rectangular se usa
\[
A=bh,\qquad I=\frac{bh^3}{12},\qquad
G=\frac{E}{2(1+\nu)},\qquad
K_s=\kappa GA,\quad\kappa=\frac56.
\]
La deformación transversal se separa en flexión $w_b$ y cortante $w_s$.
El signo negativo denota descenso. Se dan cifras suficientes para auditoría,
sin atribuir significado físico a los últimos decimales de una solución en
punto flotante.

\begin{table}[H]\centering
\caption{Resumen de los cinco casos de referencia}
\begin{tabular}{@{}llll@{}}\toprule
ID & Sistema & Acción & Comprobación principal\\\midrule
EX-01 & Voladizo horizontal & $P=10\kN$ & Flecha de punta y empotramiento\\
EX-02 & Viga biapoyada & $q=10\kN/\m$ & Reacciones, $V$, $M$, flecha\\
EX-03 & Viga biempotrada & $q=10\kN/\m$ & Momentos de empotramiento\\
EX-04 & Voladizo inclinado & $F_x=4$, $F_y=-5\kN$ & Transformación local/global\\
EX-05 & Pórtico de un vano & $F_x=5+7\kN$ & Compatibilidad y reacciones\\
\bottomrule\end{tabular}
\end{table}

\section{Datos comunes y trazabilidad}
Para EX-01, EX-04 y EX-05: $E=30000\MPa$, $\nu=0.20$,
$b=0.20\m$, $h=0.30\m$, $A=0.0600\mathrm{m}^2$,
$I=0.000450\mathrm{m}^4$, $EI=13500\kNm^2$ y
$K_s=625000\kN$.
Para EX-02 y EX-03: $E=25000\MPa$, $\nu=0.20$,
$b=0.30\m$, $h=0.50\m$, $A=0.1500\mathrm{m}^2$,
$I=0.003125\mathrm{m}^4$, $EI=78125\kNm^2$ y
$K_s=1302083.333\kN$. No se agregan rótulas ni desplazamientos impuestos.

El insumo y los resultados completos, incluidos los esfuerzos de extremo,
se generan con \texttt{bun run test:engineering-review-examples} en
\texttt{test-results/engineering-review-examples.json}. El archivo WASM
utilizado tiene SHA-256
\texttt{c9f5811354f513f8ea27ae8860e484a1e60c2984fbe5e03f9c4a16ed6c17c6d6}.
La instantánea de trabajo no es una publicación inmutable; antes de emitir
un dictamen se debe fijar un commit, un binario y los hashes de todos los
insumos.

\caseheading{EX-01}{Voladizo con carga puntual en el extremo}
\begin{figure}[H]\centering
\begin{tikzpicture}[x=1cm,y=1cm]
\fill[pattern=north east lines] (-.38,-.65) rectangle (-.18,.65);
\draw[very thick] (-.18,-.65)--(-.18,.65);
\draw[member] (-.18,0)--(7,0);
\draw[force] (7,1.2)--(7,.12) node[pos=0,above] {$P=10\kN$};
\draw[<->] (0,-.55)--(7,-.55) node[midway,fill=white] {$L=3\m$};
\node[below] at (0,-.7) {$A$}; \node[below] at (7,-.7) {$B$};
\draw[defshape] (0,-.04).. controls (2,-.08) and (5,-.35)..(7,-.63);
\node[armblue,right] at (7,-.63) {$w_B$ (amplificada)};
\end{tikzpicture}
\caption{Modelo y deformada esquemática; la deformación no está a escala.}
\end{figure}
Con $L=3\m$ y $P=10\kN$, las ecuaciones independientes son
\[
R_{Ay}=P=10\kN,\qquad |M_A|=PL=30\kNm,
\]
\[
w_B=-\frac{PL^3}{3EI}-\frac{PL}{K_s}
=-6.666667\mm-0.048000\mm=-6.714667\mm.
\]
Si $x$ se mide desde el empotramiento, el diagrama convencional tiene
$V(x)=-10\kN$ y $M(x)=-10(3-x)\kNm$. La rotación física de punta
es $-PL^2/(2EI)=-3.333333$ mrad; ARM la muestra como
$+3.333333$ mrad por su convención horaria.
\begin{figure}[H]\centering
\begin{tikzpicture}[x=1cm,y=1cm]
\node[left] at (0,1.05) {$V$}; \draw[gray] (0,1.05)--(6,1.05);
\draw[armred,thick] (0,.48)--(6,.48);
\draw[armred] (0,1.05)--(0,.48) (6,1.05)--(6,.48);
\node[right] at (6,.48) {$-10\kN$};
\node[left] at (0,-.4) {$M$}; \draw[gray] (0,-.4)--(6,-.4);
\draw[armblue,thick] (0,-1.45)--(6,-.4);
\draw[armblue] (0,-.4)--(0,-1.45);
\node[below] at (0,-1.45) {$-30\kNm$};
\node[below] at (6,-.45) {$0$};
\end{tikzpicture}
\caption{Cortante y momento de referencia, convención de cara izquierda.}
\end{figure}
\begin{table}[H]\centering
\caption{EX-01: contraste numérico}
\begin{tabular}{@{}lrrr@{}}\toprule
Magnitud & Referencia & TS & Rust\\\midrule
$R_{Ay}$ (kN) & 10.000000 & 10.000000 & 10.000000\\
$|M_A|$ (kN m) & 30.000000 & 30.000000 & 30.000000\\
$u_{y,B}$ (mm) & -6.714667 & -6.714667 & -6.714667\\
\bottomrule\end{tabular}
\end{table}

\caseheading{EX-02}{Viga simplemente apoyada con carga uniforme}
\begin{figure}[H]\centering
\begin{tikzpicture}[x=1cm,y=1cm]
\draw[member] (0,0)--(8,0);
\foreach \x in {.4,1.2,2,2.8,3.6,4.4,5.2,6,6.8,7.6}
  \draw[force] (\x,.75)--(\x,.08);
\draw[armred] (.4,.77)--(7.6,.77);
\node[above] at (4,.77) {$q=10\kN/\m$};
\draw[thick] (0,0)--(-.28,-.42)--(.28,-.42)--cycle;
\draw[thick] (-.42,-.45)--(.42,-.45);
\draw[thick] (8,0)--(7.72,-.42)--(8.28,-.42)--cycle;
\draw[thick] (7.55,-.48)--(8.45,-.48);
\draw[<->] (0,-.9)--(8,-.9) node[midway,fill=white] {$L=6\m$};
\node[below] at (0,-.9) {$A$}; \node[below] at (8,-.9) {$B$};
\end{tikzpicture}
\caption{Viga con apoyo articulado en $A$ y rodillo en $B$.}
\end{figure}
\[
R_{Ay}=R_{By}=\frac{qL}{2}=30\kN,\quad
V(x)=30-10x,\quad M(x)=30x-5x^2.
\]
\[
M(L/2)=\frac{qL^2}{8}=45\kNm,\qquad
w(L/2)=-\frac{5qL^4}{384EI}-\frac{qL^2}{8K_s}
=-2.160000-0.034560=-2.194560\mm.
\]
\begin{figure}[H]\centering
\begin{tikzpicture}[x=1cm,y=1cm]
\node[left] at (0,1.2) {$V$}; \draw[gray] (0,1.2)--(8,1.2);
\draw[armred,thick] (0,2.05)--(8,.35);
\draw[densely dotted] (4,.2)--(4,2.5);
\node[above] at (0,2.05) {$+30$}; \node[below] at (8,.35) {$-30\kN$};
\node[left] at (0,-.35) {$M$}; \draw[gray] (0,-.35)--(8,-.35);
\draw[armblue,thick,domain=0:8,samples=41]
  plot (\x,{-.35+1.3*(\x/4)*(2-\x/4)});
\node[above] at (4,.95) {$+45\kNm$};
\node[below] at (0,-.35) {$0$}; \node[below] at (8,-.35) {$6\m$};
\end{tikzpicture}
\caption{Diagramas de referencia; la abscisa representa $x=0$ a $6$ m.}
\end{figure}
\begin{table}[H]\centering
\caption{EX-02: valores para reconstruir los diagramas}
\begin{tabular}{@{}lrrrrr@{}}\toprule
$x$ (m) & 0 & 1.5 & 3.0 & 4.5 & 6.0\\\midrule
$V$ (kN) & 30 & 15 & 0 & -15 & -30\\
$M$ (kN m) & 0 & 33.75 & 45 & 33.75 & 0\\
\bottomrule\end{tabular}
\end{table}
\begin{table}[H]\centering
\caption{EX-02: contraste numérico}
\begin{tabular}{@{}lrrr@{}}\toprule
Magnitud & Referencia & TS & Rust\\\midrule
$R_{Ay}=R_{By}$ (kN) & 30.000000 & 30.000000 & 30.000000\\
$|M_{L/2}|$ (kN m) & 45.000000 & 45.000000 & 45.000000\\
$u_{y,L/2}$ (mm) & -2.194560 & -2.194560 & -2.194560\\
\bottomrule\end{tabular}
\end{table}
\note{ARM presenta $-45$ kN m en el campo de diagrama central. Su magnitud
coincide con el momento sagante $+45$ kN m adoptado aquí; para auditar signos
se debe usar una convención única en toda la comparación.}

\caseheading{EX-03}{Viga biempotrada con carga uniforme}
Mismas propiedades, luz y carga que EX-02, pero con traslaciones y giros
restringidos en ambos extremos. La solución independiente es
\[
R_{Ay}=R_{By}=30\kN,\quad
M(x)=-\frac{qL^2}{12}+\frac{qL}{2}x-\frac q2 x^2
=-30+30x-5x^2\quad(\kNm),
\]
\[
M_A=M_B=-30\kNm,\qquad M_{L/2}=+15\kNm,
\]
\[
w(L/2)=-\frac{qL^4}{384EI}-\frac{qL^2}{8K_s}
=-0.432000-0.034560=-0.466560\mm.
\]
\begin{figure}[H]\centering
\begin{tikzpicture}[x=1cm,y=1cm]
\fill[pattern=north east lines] (-.35,-.55) rectangle (-.18,.55);
\fill[pattern=north east lines] (8.18,-.55) rectangle (8.35,.55);
\draw[very thick] (-.18,-.55)--(-.18,.55) (8.18,-.55)--(8.18,.55);
\draw[member] (-.18,0)--(8.18,0);
\foreach \x in {.4,1.2,2,2.8,3.6,4.4,5.2,6,6.8,7.6}
  \draw[force] (\x,.75)--(\x,.08);
\node[above] at (4,.8) {$q=10\kN/\m$};
\draw[gray] (0,-1.15)--(8,-1.15);
\draw[armblue,thick,domain=0:8,samples=41]
  plot (\x,{-1.15+(-30+30*(\x*6/8)-5*(\x*6/8)^2)/35});
\node[below] at (0,-2) {$-30$};
\node[above] at (4,-.72) {$+15\kNm$};
\node[below] at (8,-2) {$-30$};
\end{tikzpicture}
\caption{Viga biempotrada y diagrama de momento de referencia.}
\end{figure}
\begin{table}[H]\centering
\caption{EX-03: valores del diagrama y contraste}
\begin{tabular}{@{}lrrrrr@{}}\toprule
$x$ (m) & 0 & 1.5 & 3.0 & 4.5 & 6.0\\\midrule
$V$ (kN) & 30 & 15 & 0 & -15 & -30\\
$M$ (kN m) & -30 & 3.75 & 15 & 3.75 & -30\\
\bottomrule\end{tabular}
\medskip
\begin{tabular}{@{}lrrr@{}}\toprule
Magnitud & Referencia & TS & Rust\\\midrule
$R_{Ay}=R_{By}$ (kN) & 30.000000 & 30.000000 & 30.000000\\
$|M_A|=|M_B|$ (kN m) & 30.000000 & 30.000000 & 30.000000\\
$|M_{L/2}|$ (kN m) & 15.000000 & 15.000000 & 15.000000\\
$u_{y,L/2}$ (mm) & -0.466560 & -0.466560 & -0.466560\\
\bottomrule\end{tabular}
\end{table}
\note{Este caso no contiene grados de libertad libres. Por ello comprueba
reacciones, cargas equivalentes y recuperación de deformada, pero no
ejercita la resolución de un sistema lineal libre; debe leerse junto con los
otros casos, no como prueba aislada del solucionador.}

\caseheading{EX-04}{Voladizo inclinado con cargas globales}
El elemento va de $A=(0,0)$ a $B=(3,4)$ m; $L=5$ m,
$c=3/5=0.6$, $s=4/5=0.8$. En $B$ actúan
$(F_x,F_y)=(+4,-5)$ kN.
\begin{figure}[H]\centering
\begin{tikzpicture}[x=1cm,y=1cm]
\fill[pattern=north east lines] (-.45,-.4) rectangle (.45,-.25);
\draw[very thick] (-.45,-.25)--(.45,-.25);
\draw[member] (0,0)--(3,4);
\fill (0,0) circle (1.4pt) (3,4) circle (1.4pt);
\draw[force] (3,4)--(4.25,4) node[right] {$F_x=+4\kN$};
\draw[force] (3,4)--(3,2.85) node[below] {$F_y=-5\kN$};
\draw[axis] (-.2,-.8)--(1,-.8) node[right] {$X$};
\draw[axis] (-.2,-.8)--(-.2,.3) node[above] {$Y$};
\node[left] at (0,0) {$A$}; \node[above left] at (3,4) {$B$};
\node[right] at (1.35,2.0) {$L=5\m$};
\end{tikzpicture}
\caption{Proyección de carga global sobre los ejes locales de la barra.}
\end{figure}
Las componentes locales son
\[
N=cF_x+sF_y=-1.6\kN,\qquad
T=-sF_x+cF_y=-6.2\kN.
\]
Con $u'=NL/(EA)$ y
$v'=T[L^3/(3EI)+L/K_s]$, la transformación a global
$u_x=cu'-sv'$, $u_y=su'+cv'$ entrega
\[
u_{x,B}=+15.345655\mm,\qquad
u_{y,B}=-11.514797\mm.
\]
La estática da $R_{Ax}=-4\kN$, $R_{Ay}=+5\kN$ y
$|M_A|=|3(-5)-4(4)|=31\kNm$.
\begin{table}[H]\centering
\caption{EX-04: contraste de respuesta global}
\begin{tabular}{@{}lrrr@{}}\toprule
Magnitud & Referencia & TS & Rust\\\midrule
$R_{Ax}$ (kN) & -4.000000 & -4.000000 & -4.000000\\
$R_{Ay}$ (kN) & +5.000000 & +5.000000 & +5.000000\\
$|M_A|$ (kN m) & 31.000000 & 31.000000 & 31.000000\\
$u_{x,B}$ (mm) & +15.345655 & +15.345655 & +15.345655\\
$u_{y,B}$ (mm) & -11.514797 & -11.514797 & -11.514797\\
\bottomrule\end{tabular}
\end{table}

\caseheading{EX-05}{Pórtico empotrado de un vano bajo carga lateral}
Nudos $P_0=(0,0)$, $P_1=(0,3)$, $P_2=(5,3)$,
$P_3=(5,0)$ m. Apoyos empotrados en $P_0$ y $P_3$.
Tres barras prismáticas continuas: $P_0P_1$, $P_1P_2$, $P_2P_3$.
Fuerzas horizontales $+5$ kN en $P_1$ y $+7$ kN en $P_2$.
\begin{figure}[H]\centering
\begin{tikzpicture}[x=1cm,y=1cm]
\draw[member] (0,0)--(0,3)--(5,3)--(5,0);
\fill[pattern=north east lines] (-.4,-.35) rectangle (.4,-.2);
\fill[pattern=north east lines] (4.6,-.35) rectangle (5.4,-.2);
\draw[very thick] (-.4,-.2)--(.4,-.2) (4.6,-.2)--(5.4,-.2);
\draw[force] (0,3)--(1.05,3) node[midway,above] {$5\kN$};
\draw[force] (5,3)--(6.05,3) node[right] {$7\kN$};
\node[left] at (0,0) {$P_0$}; \node[left] at (0,3) {$P_1$};
\node[above] at (5,3) {$P_2$}; \node[right] at (5,0) {$P_3$};
\draw[<->] (0,3.6)--(5,3.6) node[midway,fill=white] {$5\m$};
\draw[<->] (-.75,0)--(-.75,3) node[midway,fill=white] {$3\m$};
\draw[defshape] (0,0)..controls (.12,1.5) and (.22,2.7)..(.28,3)
  --(5.28,3)..controls (5.22,2.4) and (5.11,.8)..(5,0);
\node[armblue,right] at (5.3,2.3) {deformada amplificada};
\end{tikzpicture}
\caption{Pórtico y deformada cualitativa, no a escala.}
\end{figure}
La referencia independiente se obtiene por flexibilidad/compatibilidad de
los redundantes en el pórtico. No se usa la matriz global del motor ARM.
Los valores de la tabla son los hitos para contrastar una solución externa.
\begin{table}[H]\centering
\caption{EX-05: desplazamientos nodales; $\theta$ positiva horaria}
\begin{tabular}{@{}llrrr@{}}\toprule
Nudo & Magnitud & Referencia & TS & Rust\\\midrule
$P_1$ & $u_x$ (mm) & 1.687063 & 1.687063 & 1.687063\\
$P_1$ & $u_y$ (mm) & 0.004681 & 0.004681 & 0.004681\\
$P_1$ & $\theta$ (mrad) & 0.439239 & 0.439239 & 0.439239\\
$P_2$ & $u_x$ (mm) & 1.689831 & 1.689831 & 1.689831\\
$P_2$ & $u_y$ (mm) & -0.004681 & -0.004681 & -0.004681\\
$P_2$ & $\theta$ (mrad) & 0.440290 & 0.440290 & 0.440290\\
\bottomrule\end{tabular}
\end{table}
\begin{table}[H]\centering
\caption{EX-05: reacciones; momentos mostrados como magnitud}
\begin{tabular}{@{}llrrr@{}}\toprule
Apoyo & Magnitud & Referencia & TS & Rust\\\midrule
$P_0$ & $R_x$ (kN) & -5.996529 & -5.996529 & -5.996529\\
$P_0$ & $R_y$ (kN) & -2.808424 & -2.808424 & -2.808424\\
$P_0$ & $|M|$ (kN m) & 10.971367 & 10.971367 & 10.971367\\
$P_3$ & $R_x$ (kN) & -6.003471 & -6.003471 & -6.003471\\
$P_3$ & $R_y$ (kN) & +2.808424 & +2.808424 & +2.808424\\
$P_3$ & $|M|$ (kN m) & 10.986514 & 10.986514 & 10.986514\\
\bottomrule\end{tabular}
\end{table}
Como control estático, $\sum F_x=12-5.996529-6.003471=0$,
$\sum F_y=-2.808424+2.808424=0$.
Respecto a $P_0$, el momento externo es $-36$ kN m, el par de
reacciones verticales es $5(2.808424)=14.042119$ kN m y los
momentos físicos de base suman $21.957881$ kN m:
la suma es aproximadamente cero. ARM muestra los momentos de
reacción como $-10.971367$ y $-10.986514$ kN m; se invierte su signo
al escribir el equilibrio físico antihorario.

\section{Lectura crítica de los resultados}
Las tablas muestran concordancia al redondeo indicado en las magnitudes
seleccionadas. En la ejecución reproducible, la máxima diferencia absoluta
informada por el comparador TS--Rust entre fuerzas, reacciones y diagramas
de EX-01 a EX-05 fue menor que $8\times10^{-13}$ en las unidades del campo
comparado; para coordenadas de deformada fue menor que
$4\times10^{-16}$ m. Estos máximos no son porcentajes ni una cota de error
físico y no prueban todos los elementos o casos de carga.

Un resultado es \emph{coherente} si conserva equilibrio, satisface
restricciones, tiene signos acordes a una convención declarada y reproduce
una referencia independiente dentro de tolerancias justificadas. La
concordancia mutua de dos implementaciones que comparten hipótesis no
garantiza corrección. Sería \emph{incongruente}, por ejemplo, que EX-02 no
sumara $60$ kN de reacción vertical, que EX-04 no reaccionara con $-4$ kN
horizontal, o que EX-05 presentara deriva opuesta a las cargas sin una causa
física adicional. Los redondeos de seis decimales no deben interpretarse
como precisión experimental.

\section{Protocolo propuesto para el comité revisor}
\begin{enumerate}
\item Replicar cada geometría, apoyo, sección, material y carga en un
programa independiente o mediante cálculo manual, sin ajustar parámetros
para obtener las salidas ARM.
\item Registrar formulación de viga usada: Euler--Bernoulli o
Timoshenko; esta última incluye $K_s=\kappa GA$. Si se usa
Euler--Bernoulli, comparar por separado el aporte de flexión.
\item Verificar por caso equilibrio global, compatibilidad, deformaciones,
fuerzas internas y signos. Reportar diferencia absoluta y relativa con
tolerancias y redondeos pactados previamente.
\item Para EX-05, revisar además fuerzas de extremo, diagramas completos,
orientación local de cada barra y sensibilidad a la rigidez axial/cortante.
\item Repetir con modelos más complejos: mallas refinadas, carga parcial,
temperatura, asentamientos, rótulas, inestabilidad y geometrías irregulares;
incluir pruebas negativas y de convergencia.
\item Emitir dictamen delimitado por versión, alcance, casos y excepciones;
no extrapolar estos ensayos a cumplimiento RNE E.020, E.030 o E.060.
\end{enumerate}

\begin{table}[H]\centering
\caption{Formato mínimo para registrar discrepancias}
\begin{tabularx}{\textwidth}{@{}lXlll@{}}\toprule
Caso & Magnitud / ubicación & Referencia & ARM & Dictamen\\\midrule
\rule{0pt}{8mm}\phantom{EX-00} & & & & \\
\rule{0pt}{8mm} & & & & \\
\rule{0pt}{8mm} & & & & \\
\bottomrule\end{tabularx}
\end{table}

\section{Fuentes y archivos de revisión}
\begin{itemize}
\item Formulación conceptual del elemento Timoshenko:
\href{https://oit.tudelft.nl/Finite-Elements-in-CEG/main/structural_linear/Exercises/pyjive_timoshenko.html}
{TU Delft, ejercicio FEM Timoshenko}; teoría de vigas:
\href{https://mechanics.ju.se/SolidMechanics/BeamTheory.html}
{Jönköping University, Beam Theory}.
\item Datos de cada ensayo: \texttt{tests/engineering-review-examples.mjs};
referencia por compatibilidad: \texttt{tests/fixtures/portal-force-reference.mjs};
salida completa: \texttt{test-results/engineering-review-examples.json}.
\item Para encargo ciego sin resultados ARM:
\texttt{docs/BAS-01\_ENCARGO\_REVISION\_INDEPENDIENTE.md}.
\end{itemize}

\note{Conclusión delimitada: los cinco ensayos seleccionados pasaron la
comparación automatizada con sus referencias bajo las hipótesis declaradas.
El software permanece en fase de verificación técnica y no debe presentarse
como certificado ni aprobado para proyectos reales por este informe.}
\end{document}
